\documentclass[11pt,a4paper]{article}
\usepackage[margin=2.5cm]{geometry}
\usepackage{amsmath,amssymb}
\usepackage{booktabs}
\usepackage{parskip}
\usepackage{hyperref}
\usepackage{xcolor}

\hypersetup{colorlinks=true,linkcolor=blue}

\title{\textbf{Note: Determining the Coherent Scatter Amplitude $\varepsilon$\\
from Real Data}\\[0.4em]
\large \texttt{coh\_amplitude} calibration beyond simulations}
\author{destacking project}
\date{2026 May 23}

\begin{document}
\maketitle

\section*{Context}

The rank-1 covariance inflation scheme adds a term
$\varepsilon^2\,\mathbf{m}_i\mathbf{m}_i^\top$ to each cluster's covariance
matrix to absorb coherent per-cluster amplitude scatter from triaxiality and
miscentering.  The calibration scan approach finds the value of $\varepsilon$
that makes the posterior width match the run-to-run scatter of the estimator
in simulations.  The question is: \emph{how would we determine $\varepsilon$
when we only have one realisation of the real data?}

\section{The problem}

In simulations we have access to many independent realisations of the survey,
so we can measure $\sigma_{\rm runs}$ directly.  In real data we have one
survey, one cluster sample, and an unknown true $\gamma$.  The simulation-based
calibration ratio $\sigma_{\rm runs}/\langle\sigma_{\rm post}\rangle$ cannot be
computed.

\section{Available approaches}

\subsection{Simulation-based transfer}
Tune $\varepsilon$ in simulations that match the real survey (cluster
redshift and richness distributions, source $n(z)$, shape noise).  If the
simulations faithfully reproduce the physical sources of coherent scatter —
triaxiality distribution, miscentering distribution — then the calibrated
$\varepsilon$ transfers to the real data.

\textbf{Risk:} if the real coherent scatter amplitude differs from what the
simulations inject, $\varepsilon$ will be wrong and error bars will be
miscalibrated.  Validation requires an independent check (e.g.\
cross-validation or comparison with external miscentering/triaxiality
measurements).

\subsection{Joint inference: treat $\varepsilon$ as a nuisance parameter}
Add $\varepsilon$ to the sampler's parameter space and infer it jointly with
$\gamma$, marginalising the posterior over $\varepsilon$ with a physically
motivated prior.

\textbf{Risk:} requires the likelihood to genuinely constrain $\varepsilon$
from the data.  With the current rank-1 parameterisation this is dangerous:
the scan has shown that the likelihood develops pathological modes at
$\varepsilon \gtrsim 0.2$.  This approach should not be attempted until the
spurious $\hat\gamma \approx +1.4$ mode is understood and eliminated.

\subsection{Empirical two-step variance estimation}
Measure the coherent amplitude variance directly from the data without
involving $\gamma$:
\begin{enumerate}
  \item Fit a mean NFW profile to the stacked cluster sample.
  \item For each cluster, measure the residual amplitude relative to
        the mean profile.
  \item The component of that residual scatter correlated with the model
        direction $\mathbf{m}_i$ gives $\varepsilon$ independently of $\gamma$.
\end{enumerate}
This decouples noise-model estimation from the scientific inference.

\textbf{Risk:} the amplitude scatter measured from noisy individual profiles
will be upward-biased; a noise-bias correction is required.  Also requires
that the mean profile is a good reference, which it may not be if the sample
has a broad mass/redshift range.

\subsection{Physical priors from external constraints}
Predict $\varepsilon$ from independent measurements rather than fitting it
from lensing data at all:
\begin{itemize}
  \item Triaxiality: halo shape distributions from N-body simulations
        (e.g.\ Jing \& Suto 2002, Allgood et al.\ 2006).
  \item Miscentering: BCG offset or X-ray centroid distributions from
        optical/X-ray data (e.g.\ Johnston et al.\ 2007, Zhang et al.\ 2019).
\end{itemize}
Both effects predict a coherent amplitude scatter of order $10$--$20\%$ per
cluster; translated to the model-direction variance, this gives a physically
grounded prior or fixed value for $\varepsilon$.

\textbf{Risk:} the coherent scatter in real data is not purely along the
model direction, so any single $\varepsilon$ is an approximation regardless
of how well the physical inputs are known.

\section{Recommendation and priority ordering}

\begin{center}
\begin{tabular}{r l l}
\toprule
Priority & Approach & Prerequisite \\
\midrule
1 & Fix the spurious $\hat\gamma \approx +1.4$ mode & (current blocker) \\
2 & Physical priors from N-body + miscentering data & none \\
3 & Simulation-based transfer & fix (1) \\
4 & Empirical two-step & fix (1), noise-bias correction \\
5 & Joint inference over $\varepsilon$ & fix (1), validate no new modes \\
\bottomrule
\end{tabular}
\end{center}

In practice the community typically combines approaches (2) and (3): simulations
are tuned to known physical scatter sources and validated where possible against
the data.  The key message is that if the simulations do not match the real
coherent scatter amplitude, the calibrated $\varepsilon$ will be wrong
regardless of how well the scan is performed — which makes fixing the
pathological mode a prerequisite for any of these approaches to be trustworthy.

\end{document}
